% chktex-file 44
\section{Port Descriptions}

The ports for \textbf{Mac} are shown below in
Table~\ref{table:ports}. The width of several ports is controlled
by the following parameters:

\begin{itemize}[noitemsep]
  \item \textit{dataW} is the width of the XGMII data bus in bits
  \item \textit{ctrlW} is the width of the XGMII control bus in bits, which must be equal to \textit{dataW}/8 for byte granularity
  \item \textit{gbxCnt} is the width of the internal counter used for tracking the number of bytes in a frame, which must be wide enough to count up to the maximum frame size (e.g., 12 bits for 4096 bytes)
  \item \textit{ptpTsW} is the width of the PTP timestamp bus in bits, which must be sufficient to represent the desired timestamp resolution and range (e.g., 48 bits for nanosecond resolution over several years)
  
\end{itemize}

\renewcommand*{\arraystretch}{1.3}
\begin{longtable}[H]{
    | p{0.30\textwidth}
    | p{0.15\textwidth}
    | p{0.10\textwidth}
    | p{0.35\textwidth} |
  }
  \hline
  \textbf{Port Name} & \textbf{Width} & \textbf{Direction} & \textbf{Description} \\ \hline \hline
  \endhead

  % --- Clock and Reset ---
  \rowcolor{fog-grey} \multicolumn{4}{|l|}{\textbf{Clock and Reset Signals}} \\ \hline
  rxClk              & 1              & Input  & RX domain clock \\ \hline
  rxRst              & 1              & Input  & RX domain active-high synchronous reset \\ \hline
  txClk              & 1              & Input  & TX domain clock \\ \hline
  txRst              & 1              & Input  & TX domain active-high synchronous reset \\ \hline
  statClk            & 1              & Input  & Statistics domain clock \\ \hline
  statRst            & 1              & Input  & Statistics domain active-high synchronous reset \\ \hline

  % --- AXI-Stream Interfaces ---
  \rowcolor{fog-grey} \multicolumn{4}{|l|}{\textbf{User Side AXI-Stream Interfaces}} \\ \hline
  sAxisTx            & See Table~\ref{table:axi_signals} & \textbf{Sink} (Slave)  & \textbf{TX Sink:} User data to be transmitted \\ \hline
  mAxisTxCpl         & See Table~\ref{table:axi_signals} & \textbf{Source} (Master) & \textbf{TX Completion:} Status and PTP timestamps \\ \hline
  mAxisRx            & See Table~\ref{table:axi_signals} & \textbf{Source} (Master) & \textbf{RX Source:} Received Ethernet frames to user \\ \hline
  mAxisStat          & See Table~\ref{table:axi_signals} & \textbf{Source} (Master) & \textbf{Stats Source:} Aggregated statistics update bus \\ \hline
  
  % --- XGMII ---
  \rowcolor{fog-grey} \multicolumn{4}{|l|}{\textbf{Physical Layer Interface (XGMII)}} \\ \hline
  xgmiiRxd           & \textit{dataW} & Input  & XGMII receive data bus \\ \hline
  xgmiiRxc           & \textit{ctrlW} & Input  & XGMII receive control bus \\ \hline
  xgmiiRxValid       & 1              & Input  & XGMII receive data valid \\ \hline
  xgmiiTxd           & \textit{dataW} & Output & XGMII transmit data bus \\ \hline
  xgmiiTxc           & \textit{ctrlW} & Output & XGMII transmit control bus \\ \hline
  xgmiiTxValid       & 1              & Output & XGMII transmit data valid \\ \hline

  % --- Gearbox ---
  \rowcolor{fog-grey} \multicolumn{4}{|l|}{\textbf{Gearbox and Synchronization}} \\ \hline
  txGbxReqSync       & \textit{gbxCnt}& Input  & Gearbox synchronization request from PCS \\ \hline
  txGbxReqStall      & 1              & Input  & Gearbox stall request (pauses MAC) \\ \hline
  txGbxSync          & \textit{gbxCnt}& Output & Gearbox synchronization acknowledgment \\ \hline

  % --- PTP ---
  \rowcolor{fog-grey} \multicolumn{4}{|l|}{\textbf{Precision Time Protocol (PTP)}} \\ \hline
  txPtpTs            & \textit{ptpTsW}& Input  & Reference timestamp for TX SOF capture \\ \hline
  rxPtpTs            & \textit{ptpTsW}& Input  & Reference timestamp for RX SOF capture \\ \hline

  % --- LFC ---
  \rowcolor{fog-grey} \multicolumn{4}{|l|}{\textbf{Link-Level Flow Control (LFC)}} \\ \hline
  txLfcReq           & 1              & Input  & Trigger LFC frame transmission \\ \hline
  txLfcResend        & 1              & Input  & Retransmit current LFC pause frame \\ \hline
  rxLfcEn            & 1              & Input  & Enable LFC frame reception/processing \\ \hline
  rxLfcReq           & 1              & Output & Asserted when an LFC pause is received \\ \hline
  rxLfcAck           & 1              & Input  & Acknowledge LFC pause request \\ \hline

  % --- PFC ---
  \rowcolor{fog-grey} \multicolumn{4}{|l|}{\textbf{Priority Flow Control (PFC)}} \\ \hline
  txPfcReq           & 8              & Input  & Trigger PFC for 8 priority classes \\ \hline
  txPfcResend        & 1              & Input  & Retransmit current PFC pause frame \\ \hline
  rxPfcEn            & 8              & Input  & Enable PFC for 8 priority classes \\ \hline
  rxPfcReq           & 8              & Output & Asserted when a PFC pause is received \\ \hline
  rxPfcAck           & 8              & Input  & Acknowledge PFC pause request \\ \hline

  % --- Pause ---
  \rowcolor{fog-grey} \multicolumn{4}{|l|}{\textbf{Pause Interface}} \\ \hline
  txLfcPauseEn       & 1              & Input  & Enable automatic TX pause on LFC \\ \hline
  txPauseReq         & 1              & Input  & Trigger standard 802.3x pause frame \\ \hline
  txPauseAck         & 1              & Output & Acknowledge pause frame transmission \\ \hline

  % --- TX Status ---
  \rowcolor{fog-grey} \multicolumn{4}{|l|}{\textbf{Transmit Status and Monitoring}} \\ \hline
  txStartPacket      & 2              & Output & Pulse indicating SOF (lane aligned) \\ \hline
  statTxByte         & 4              & Output & Valid bytes in current TX cycle \\ \hline
  statTxPktLen       & 16             & Output & Length of last transmitted packet \\ \hline
  statTxPktUcast     & 1              & Output & Last TX packet was Unicast \\ \hline
  statTxPktMcast     & 1              & Output & Last TX packet was Multicast \\ \hline
  statTxPktBcast     & 1              & Output & Last TX packet was Broadcast \\ \hline
  statTxPktVlan      & 1              & Output & Last TX packet was VLAN tagged \\ \hline
  statTxPktGood      & 1              & Output & Last TX packet was valid \\ \hline
  statTxPktBad       & 1              & Output & Last TX packet had errors \\ \hline
  statTxErrOversize  & 1              & Output & TX Packet exceeded MTU \\ \hline
  statTxErrUser      & 1              & Output & TX Error triggered by user (tuser) \\ \hline
  statTxErrUnderflow & 1              & Output & TX Error: sAxisTx underflowed \\ \hline

  % --- RX Status ---
  \rowcolor{fog-grey} \multicolumn{4}{|l|}{\textbf{Receive Status and Monitoring}} \\ \hline
  rxStartPacket      & 2              & Output & Pulse indicating SOF (lane aligned) \\ \hline
  statRxByte         & 4              & Output & Valid bytes in current RX cycle \\ \hline
  statRxPktLen       & 16             & Output & Length of last received packet \\ \hline
  statRxPktFragment  & 1              & Output & Received packet was a fragment ($<$64B) \\ \hline
  statRxPktJabber    & 1              & Output & Received packet was a jabber ($>$MTU) \\ \hline
  statRxPktUcast     & 1              & Output & Last RX packet was Unicast \\ \hline
  statRxPktMcast     & 1              & Output & Last RX packet was Multicast \\ \hline
  statRxPktBcast     & 1              & Output & Last RX packet was Broadcast \\ \hline
  statRxPktVlan      & 1              & Output & Last RX packet was VLAN tagged \\ \hline
  statRxPktGood      & 1              & Output & Last RX packet passed FCS check \\ \hline
  statRxPktBad       & 1              & Output & Last RX packet failed FCS check \\ \hline
  statRxErrOversize  & 1              & Output & RX Packet exceeded MTU \\ \hline
  statRxErrBadFcs    & 1              & Output & RX Packet had a bad FCS residue \\ \hline
  statRxErrBadBlock  & 1              & Output & RX Packet contained 64b/66b error blocks \\ \hline
  statRxErrFraming   & 1              & Output & RX Packet had XGMII framing errors \\ \hline
  statRxErrPreamble  & 1              & Output & RX Packet had an invalid preamble \\ \hline
  statRxFifoDrop     & 1              & Input  & External RX FIFO overflow indicator \\ \hline

  % --- Control Statistics ---
  \rowcolor{fog-grey} \multicolumn{4}{|l|}{\textbf{Control Frame Statistics}} \\ \hline
  statTxMcf          & 1              & Output & Transmitted MAC Control Frame \\ \hline
  statRxMcf          & 1              & Output & Received MAC Control Frame \\ \hline
  statTxLfcPkt       & 1              & Output & Transmitted LFC packet \\ \hline
  statTxLfcXon       & 1              & Output & Transmitted LFC XON (Pause 0) \\ \hline
  statTxLfcXoff      & 1              & Output & Transmitted LFC XOFF \\ \hline
  statTxLfcPaused    & 1              & Output & TX path is currently paused by LFC \\ \hline
  statTxPfcPkt       & 1              & Output & Transmitted PFC packet \\ \hline
  statTxPfcXon       & 8              & Output & Transmitted PFC XON per class \\ \hline
  statTxPfcXoff      & 8              & Output & Transmitted PFC XOFF per class \\ \hline
  statTxPfcPaused    & 8              & Output & TX path currently paused per PFC class \\ \hline
  statRxLfcPkt       & 1              & Output & Received LFC packet \\ \hline
  statRxLfcXon       & 1              & Output & Received LFC XON \\ \hline
  statRxLfcXoff      & 1              & Output & Received LFC XOFF \\ \hline
  statRxLfcPaused    & 1              & Output & RX path reports LFC pause state \\ \hline
  statRxPfcPkt       & 1              & Output & Received PFC packet \\ \hline
  statRxPfcXon       & 8              & Output & Received PFC XON per class \\ \hline
  statRxPfcXoff      & 8              & Output & Received PFC XOFF per class \\ \hline
  statRxPfcPaused    & 8              & Output & RX path reports PFC pause per class \\ \hline

  % --- Runtime Configuration ---
  \rowcolor{fog-grey} \multicolumn{4}{|l|}{\textbf{Runtime Configuration}} \\ \hline
  cfgTxMaxPktLen     & 16             & Input  & Maximum frame length for TX \\ \hline
  cfgTxIfg           & 8              & Input  & Inter-Frame Gap (IFG) setting in bytes \\ \hline
  cfgTxEnable        & 1              & Input  & Global TX path enable \\ \hline
  cfgRxMaxPktLen     & 16             & Input  & Maximum frame length for RX \\ \hline
  cfgRxEnable        & 1              & Input  & Global RX path enable \\ \hline
  cfgMcfRxEthDstMcast     & 48        & Input  & MCF Multicast Destination MAC \\ \hline
  cfgMcfRxCheckEthDstMcast& 1         & Input  & Enable MCF Multicast address check \\ \hline
  cfgMcfRxEthDstUcast     & 48        & Input  & MCF Unicast Destination MAC \\ \hline
  cfgMcfRxCheckEthDstUcast& 1         & Input  & Enable MCF Unicast address check \\ \hline
  cfgMcfRxEthSrc          & 48        & Input  & MCF Source MAC address \\ \hline
  cfgMcfRxCheckEthSrc     & 1         & Input  & Enable MCF Source address check \\ \hline
  cfgMcfRxEthType         & 16        & Input  & MCF Ethernet Type (usually 0x8808) \\ \hline
  cfgMcfRxOpcodeLfc        & 16        & Input  & Opcode for LFC frames \\ \hline
  cfgMcfRxCheckOpcodeLfc   & 1         & Input  & Enable LFC Opcode checking \\ \hline
  cfgMcfRxOpcodePfc        & 16        & Input  & Opcode for PFC frames \\ \hline
  cfgMcfRxCheckOpcodePfc   & 1         & Input  & Enable PFC Opcode checking \\ \hline
  cfgMcfRxForward          & 1         & Input  & Forward MCF frames to user logic \\ \hline
  cfgMcfRxEnable           & 1         & Input  & Global MCF processing enable \\ \hline
  cfgTxLfcEthDst          & 48        & Input  & TX LFC Destination MAC \\ \hline
  cfgTxLfcEthSrc          & 48        & Input  & TX LFC Source MAC \\ \hline
  cfgTxLfcEthType         & 16        & Input  & TX LFC Ethernet Type \\ \hline
  cfgTxLfcOpcode          & 16        & Input  & TX LFC Opcode \\ \hline
  cfgTxLfcEn              & 1         & Input  & Enable TX LFC generation \\ \hline
  cfgTxLfcQuanta          & 16        & Input  & Pause duration in quanta for LFC \\ \hline
  cfgTxLfcRefresh         & 16        & Input  & Refresh timer value for LFC \\ \hline
  cfgTxPfcEthDst          & 48        & Input  & TX PFC Destination MAC \\ \hline
  cfgTxPfcEthSrc          & 48        & Input  & TX PFC Source MAC \\ \hline
  cfgTxPfcEthType         & 16        & Input  & TX PFC Ethernet Type \\ \hline
  cfgTxPfcOpcode          & 16        & Input  & TX PFC Opcode \\ \hline
  cfgTxPfcEn              & 1         & Input  & Enable TX PFC generation \\ \hline
  cfgTxPfcQuanta          & 8 $\times$ 16 & Input  & Pause duration quanta per priority class \\ \hline
  cfgTxPfcRefresh         & 8 $\times$ 16 & Input  & Refresh timer per priority class \\ \hline
  cfgRxLfcOpcode          & 16        & Input  & Expected Opcode for incoming LFC \\ \hline
  cfgRxLfcEn              & 1         & Input  & Enable incoming LFC processing \\ \hline
  cfgRxPfcOpcode          & 16        & Input  & Expected Opcode for incoming PFC \\ \hline
  cfgRxPfcEn              & 1         & Input  & Enable incoming PFC processing \\ \hline

  \caption{MAC Port Descriptions}\label{table:ports}
\end{longtable}


\subsection{AXI-Stream Interface Signals}
All AXI-Stream interfaces in the \textbf{Mac} core follow the signal definitions shown in Table~\ref{table:axi_signals}. 
Note that the width of some signals are determined by the specific \textit{AxisInterfaceParams} configuration for that port.

\renewcommand*{\arraystretch}{1.4}
\begin{longtable}[H]{
    | p{0.20\textwidth}
    | p{0.20\textwidth}
    | p{0.15\textwidth}
    | p{0.38\textwidth} |
  }
  \hline
  \textbf{Signal} & \textbf{Width} & \textbf{Direction\textsuperscript{1}} & \textbf{Description} \\ \hline \hline
  tdata  & \textit{dataW}  & Output & Primary data bus \\ \hline
  tkeep* & \textit{keepW}  & Output & Byte qualifiers; indicates valid bytes in \textit{tdata} \\ \hline
  tstrb* & \textit{keepW}  & Output & Position byte qualifiers (usually matches \textit{tkeep}) \\ \hline
  tid* & \textit{idW}    & Output & Stream identifier (used for TX tags/completion) \\ \hline
  tdest* & \textit{destW}  & Output & Routing destination identifier \\ \hline
  tuser* & \textit{userW}  & Output & Sideband metadata (used for PTP timestamps in RX) \\ \hline
  tlast  & 1               & Output & Indicates the last transfer of a frame \\ \hline
  tvalid & 1               & Output & Indicates that the source is driving valid signals \\ \hline
  tready & 1               & Input  & Indicates that the sink can accept data \\ \hline
  \caption{AXI-Stream Signal Definitions}\label{table:axi_signals}
\end{longtable}
\small{\textsuperscript{1} Directions are shown for a \textbf{Source} (Master) interface. For \textbf{Sink} (Slave) interfaces, such as \textit{sAxisTx}, all directions are reversed (Flipped).}


\subsection{Interface Configuration Summary}
Table~\ref{table:interface_configs} summarizes the parameter-driven widths for each AXI-Stream interface instance.

\begin{table}[H]
\centering
\begin{tabular}{|l|c|c|c|c|l|}
\hline
\textbf{Interface} & \textbf{dataW} & \textbf{keepW} & \textbf{idW} & \textbf{userW} & \textbf{destW} \\ \hline
sAxisTx    & \textit{p.dataW} & \textit{p.dataW / 8} & 8 & 1 & 8 \\ \hline
mAxisTxCpl & 96 & 1 & 8 & 1 & 8 \\ \hline
mAxisRx    & \textit{p.dataW} & \textit{p.dataW / 8} & 8 & \textit{rxUserW}\textsuperscript{2} & 8 \\ \hline
mAxisStat  & 16 & 1 & 8 & 1 & 8 \\ \hline
\end{tabular}
\caption{Interface Parameter Mapping}\label{table:interface_configs}
\small{\textsuperscript{2} \textit{rxUserW} is calculated as \textit{p.ptpTsW + 1} if PTP is enabled.}
\end{table}